\documentclass[11pt]{article}
\usepackage[a4paper,margin=27mm]{geometry}
\usepackage{amsmath,amssymb,amsthm,mathtools}
\usepackage[T1]{fontenc}
\usepackage{lmodern,microtype}
\usepackage[colorlinks=true,linkcolor=blue,urlcolor=blue]{hyperref}
\newcommand{\F}{\mathbb F}
\newcommand{\PP}{\mathbb P}
\newcommand{\A}{\mathbb A}
\newcommand{\cO}{\mathcal O}
\newcommand{\fm}{\mathfrak m}
\newcommand{\gr}{\operatorname{gr}}
\newcommand{\Nik}{\operatorname{Nik}}
\newtheorem{theorem}{Theorem}
\newtheorem{lemma}[theorem]{Lemma}
\newtheorem{proposition}[theorem]{Proposition}
\newtheorem{corollary}[theorem]{Corollary}
\theoremstyle{remark}
\newtheorem{remark}[theorem]{Remark}
\title{Nikodym sets: a power saving\\from finite-grid interpolation}
\author{Research note}
\date{5 September 2026}
\begin{document}
\maketitle

\begin{abstract}
We prove
\[
 \Nik(d,q)\ge q^d-(8d^2+1)q^{\,d-2^{1-d}}
 \qquad(d\ge2)
\]
for every prime power $q$. The main observation compares interpolation
on the whole finite grid with interpolation on an irreducible carrier.
Removing the conditions at private points supplies a polynomial cut
whose degree is independent of the degrees of the carrier's defining
equations. A dimension induction then halves the saving at each step.
The proof uses monomial counting, finite-grid Hermite interpolation,
linear Noether normalization, and the elementary degree bound for a
proper hypersurface section. It uses no resolution of singularities,
intersection pairing, or restriction on the characteristic.
\end{abstract}

\section{The statement for families with private points}

Fix an algebraic closure $K$ of $\F_q$. An affine $\F_q$-line
means its extension to a line over $K$ as well as its $q$ rational
points, as appropriate. A family $\mathcal L$ of distinct affine
$\F_q$-lines has \emph{private points} if for each $\ell\in\mathcal L$
there is a specified $b_\ell\in\ell(\F_q)$ lying on no other line of
$\mathcal L$. In particular, the $b_\ell$ are distinct.

\begin{theorem}\label{thm:carrier}
Let $d\ge2$, $1\le k\le d$, and let $A\subseteq\PP^d_K$ be an
integral projective variety of dimension $k$ and degree $\Delta$.
Suppose $A$ contains the projective closures of $L$ affine
$\F_q$-lines having private points. Then
\begin{equation}\label{eq:carrier}
 L\le C_d\Delta q^{\,k-2^{1-k}},
 \qquad C_d=8d^2+1.
\end{equation}
The variety $A$ need not be defined over $\F_q$.
\end{theorem}

\begin{corollary}\label{cor:nikodym}
For every $d\ge2$, every prime power $q$, and every Nikodym set
$N\subseteq\F_q^d$,
\begin{equation}\label{eq:main}
 |\F_q^d\setminus N|\le (8d^2+1)q^{\,d-2^{1-d}}.
\end{equation}
Consequently,
\[
 \Nik(d,q)=q^d-O_d\bigl(q^{\,d-2^{1-d}}\bigr).
\]
\end{corollary}

\begin{proof}[Deduction of the corollary]
Put $B=\F_q^d\setminus N$. For each $b\in B$, choose a witness
line $\ell_b$ with $B\cap\ell_b=\{b\}$. These lines are distinct,
and $b$ is a private point of $\ell_b$. Apply
Theorem~\ref{thm:carrier} with $A=\PP^d$, $k=d$, and $\Delta=1$.
\end{proof}

\section{Two polynomial dimension estimates}

For a projective variety $A\subseteq\PP^d$, write $H_A(t)$ for its
homogeneous Hilbert function: the dimension of degree-$t$ forms
modulo those vanishing on $A$. If $A$ meets the standard affine
chart, this is also the dimension of restrictions to $A\cap\A^d$
of affine polynomials of total degree at most $t$.

\begin{lemma}[Normalized Hilbert functions]\label{lem:normalized}
For any homogeneous ideal in $K[X_0,\ldots,X_d]$, the Hilbert
function $H$ of its quotient satisfies
\begin{equation}\label{eq:normalized}
 \frac{H(u)}{\binom{u+d}{d}}
 \le \frac{H(t)}{\binom{t+d}{d}}
 \qquad(0\le t\le u).
\end{equation}
\end{lemma}

\begin{proof}
Replace the ideal by an initial monomial ideal, which preserves its
Hilbert function. Call a monomial \emph{standard} if it is outside
this ideal. Every divisor of a standard monomial is standard.

Give the edge from $X^\alpha$ of degree $t$ to
$X^{\alpha+e_i}$ of degree $t+1$ weight $\alpha_i+1$.
Every standard child has only standard parents, and its total
incoming weight is $t+1$. The total outgoing weight from any
standard parent is $t+d+1$. Counting weighted edges therefore gives
\[
 (t+1)H(t+1)\le(t+d+1)H(t).
\]
This is the adjacent-degree case of \eqref{eq:normalized}; iterate.
\end{proof}

\begin{lemma}[Degree upper bound]\label{lem:upper}
If $A\subseteq\PP^d_K$ is integral of dimension $k$ and degree
$\Delta$, then
\begin{equation}\label{eq:upper}
 H_A(t)\le\Delta\binom{t+k}{k}\qquad(t\ge0).
\end{equation}
\end{lemma}

\begin{proof}
Here is a proof of this standard degree estimate. Choose a finite
linear projection $\pi:A\to\PP^k$ of degree $\Delta$ whose function
field extension is separable. Such a projection exists over the
infinite perfect field $K$: linear Noether normalization gives
finiteness, and generic linear coordinates are separating
coordinates on the smooth locus.

Choose an affine base point with $\Delta$ distinct unramified
preimages. The completed local rings at these preimages have the
same $k$ base coordinates. Taking the Taylor coefficients of total
degree at most $t$ at all these branches defines a linear map
from degree-$t$ polynomial restrictions into a vector space of
dimension $\Delta\binom{t+k}{k}$.

This map is injective. Indeed, if a restriction $P$ has all these
coefficients zero, its field norm to the base vanishes to order at
least $\Delta(t+1)$. On the other hand, that norm is a polynomial
of degree at most $\Delta t$. To justify the degree statement,
the homogeneous coordinate ring of $A$ is finite over the
polynomial ring in the $k+1$ projection coordinates. The norm of
a degree-$t$ homogeneous element lies in that polynomial ring
by integrality and its integral closedness, and has homogeneous
degree $\Delta t$. Dehomogenize on the chosen affine chart.

A nonzero polynomial cannot vanish at a point to order greater
than its total degree. Thus the norm is zero, which forces
$P=0$ because a nonzero field element has nonzero norm.
The Taylor map is injective, proving the bound. All Taylor
coefficients are formal coefficients; no division by factorials
is involved.
\end{proof}

\section{Interpolation on the entire finite grid}

For $x\in A\cap\A^d(K)$ and $r\ge1$, put
\[
 j_{A,x}(r)=\dim_K\bigl(\cO_{A,x}/\fm_{A,x}^{\,r}\bigr).
\]
Set $j_{A,x}(r)=0$ when $x\notin A$.

\begin{lemma}[Local minimum]\label{lem:local}
For every point $x$ of an integral $k$-dimensional variety,
\begin{equation}\label{eq:local-min}
 j_{A,x}(r)\ge\binom{r+k-1}{k}.
\end{equation}
\end{lemma}

\begin{proof}
The associated graded ring $\gr_{\fm}\cO_{A,x}$ is a standard
graded $K$-algebra of dimension $k$. Linear Noether normalization
gives $k$ algebraically independent degree-one elements in this
ring. Their monomials of degree less than $r$ are linearly
independent. There are $\binom{r+k-1}{k}$ such monomials, and
the sum of the first $r$ graded dimensions is $j_{A,x}(r)$.
\end{proof}

\begin{lemma}[Grid interpolation]\label{lem:grid}
For every integer $r\ge1$, put
\begin{equation}\label{eq:U}
 U=q(r-1)+d(q-1).
\end{equation}
Restriction of degree-at-most-$U$ polynomials on $A\cap\A^d$
surjects onto
\[
 \bigoplus_{x\in\F_q^d\cap A}
 \cO_{A,x}/\fm_{A,x}^{\,r}.
\]
In particular,
\begin{equation}\label{eq:grid-dim}
 \sum_{x\in\F_q^d}j_{A,x}(r)\le H_A(U).
\end{equation}
\end{lemma}

\begin{proof}
In $R=K[X_1,\ldots,X_d]$, let
\[
 Z_i=X_i^q-X_i,\qquad J=(Z_1,\ldots,Z_d).
\]
The polynomials $X_i^q-X_i$ have distinct roots, exactly $\F_q$.
Thus $J=\bigcap_{x\in\F_q^d}\fm_x$. The distinct maximal ideals
are pairwise comaximal, so
\[
 J^r=\bigcap_{x\in\F_q^d}\fm_x^r,\qquad
 R/J^r\cong\bigoplus_{x\in\F_q^d}R/\fm_x^r.
\]

Repeated use of $X_i^q=Z_i+X_i$ shows that $R$ is spanned over
$K[Z_1,\ldots,Z_d]$ by the monomials $X^\alpha$ with
$0\le\alpha_i<q$. Hence $R/J^r$ is spanned by
\[
 X^\alpha Z^\gamma,
 \qquad 0\le\alpha_i<q,\quad |\gamma|<r.
\]
Every such polynomial has degree at most
$d(q-1)+q(r-1)=U$. Consequently ambient polynomials of degree at
most $U$ interpolate arbitrary ambient jets of order less than $r$
at every grid point simultaneously.

Now quotient each jet space by the affine ideal of $A$. At a
point $x\in A$ the resulting quotient is precisely
$\cO_{A,x}/\fm_{A,x}^r$; at a point outside $A$ it is zero.
The map factors through polynomial restrictions on $A$.
This proves both surjectivity and \eqref{eq:grid-dim}.
\end{proof}

One useful consequence, requiring no separate point-counting theorem, is
\begin{equation}\label{eq:point-count}
 |A\cap\F_q^d|\le\Delta q^k.
\end{equation}
Indeed, Lemmas~\ref{lem:local}, \ref{lem:grid}, and \ref{lem:upper}
give
\[
 |A\cap\F_q^d|\binom{r+k-1}{k}
 \le\Delta\binom{q(r-1)+d(q-1)+k}{k}.
\]
Divide by the binomial on the left and let the integer $r$ tend
to infinity, keeping $q$ fixed.

\section{A polynomial cut independent of the defining degrees}

\begin{proposition}\label{prop:cut}
Suppose $A$ and its $L>0$ lines satisfy the hypotheses of
Theorem~\ref{thm:carrier}. Define
\[
 \rho=\frac{L}{\Delta q^k},\qquad
 a=8d^2,\qquad r=\left\lceil\frac a\rho\right\rceil.
\]
If $r\le q$, there is a polynomial $G$ of degree at most
\begin{equation}\label{eq:T}
 T=r(q-1)-1
\end{equation}
which does not vanish identically on $A$ and vanishes identically
on every selected line.
\end{proposition}

\begin{proof}
Let $B_0$ be the set of the $L$ private points. Impose on
polynomial restrictions on $A$ the linear conditions
\[
 G\in\fm_{A,x}^{\,r}
 \qquad\bigl(x\in(A\cap\F_q^d)\setminus B_0\bigr).
\]
The number of conditions is at most
\begin{align}
 \sum_{x\in\F_q^d\setminus B_0}j_{A,x}(r)
 &=\sum_{x\in\F_q^d}j_{A,x}(r)
     -\sum_{b\in B_0}j_{A,b}(r)\notag\\
 &\le H_A(U)-L\binom{r+k-1}{k},\label{eq:saved}
\end{align}
where $U$ is given by \eqref{eq:U}.

We show that the last expression is strictly less than $H_A(T)$.
By \eqref{eq:point-count}, $0<\rho\le1$, so $r\ge8d^2$.
Writing $s=U-T$, we have
\[
 s=r+(d-1)q-d+1\le dq,
 \qquad T+1=r(q-1)\ge rq/2.
\]
The normalized Hilbert inequality gives
\[
 H_A(U)-H_A(T)
 \le H_A(T)\left(
       \frac{\binom{U+d}{d}}{\binom{T+d}{d}}-1\right).
\]
Moreover,
\begin{align*}
 \frac{\binom{U+d}{d}}{\binom{T+d}{d}}
 &=\prod_{i=1}^d\left(1+\frac{s}{T+i}\right)
 \le\exp\!\left(\frac{ds}{T+1}\right)
 \le\exp\!\left(\frac{2d^2}{r}\right).
\end{align*}
Since $2d^2/r\le1/4$, the inequality $e^z-1\le2z$ for
$0\le z\le1/4$ shows that the parenthesis is at most $4d^2/r$.
Finally, for $1\le i\le k$,
\[
 T+i=r(q-1)+i-1\le q(r+i-1),
\]
and therefore
\[
 \binom{T+k}{k}\le q^k\binom{r+k-1}{k}.
\]
Combining these estimates with Lemma~\ref{lem:upper} yields
\begin{align}
 H_A(U)-H_A(T)
 &\le\frac{4d^2}{r}\Delta q^k\binom{r+k-1}{k}\notag\\
 &\le\frac{\rho}{2}\Delta q^k\binom{r+k-1}{k}
 =\frac L2\binom{r+k-1}{k}.
 \label{eq:gap}
\end{align}
Equations~\eqref{eq:saved} and \eqref{eq:gap} prove the strict
dimension inequality. There is consequently a nonzero restriction
on $A$ of a polynomial $G$ of degree at most $T$ satisfying all
the imposed conditions.

Each selected line contains exactly one point of $B_0$, its own
private point. At each of its other $q-1$ rational points,
restriction from $A$ to the line carries
$\fm_{A,x}^r$ into the $r$th power of the line's maximal ideal.
Thus $G$ restricted to the line has $q-1$ roots of multiplicity
at least $r$. Its degree is at most $T<r(q-1)$, so it vanishes
identically on the line.
\end{proof}

\section{Dimension induction and the exponent}

We recall the elementary degree fact used here. If an integral
$k$-dimensional projective variety of degree $\Delta$ is not
contained in a degree-$e$ hypersurface, its intersection has
pure dimension $k-1$, and the sum of the degrees of its reduced
irreducible components is at most $e\Delta$. The degree statement
follows by taking the Hilbert function difference
$H_A(t)-H_A(t-e)$; its leading term is
$e\Delta t^{k-1}/(k-1)!$. The principal ideal theorem and the
dimension formula give the dimension assertion. Intersection
multiplicities are positive, so discarding them can only decrease
the degree sum.

\begin{proof}[Proof of Theorem~\ref{thm:carrier}]
Fix $d$, put $a=8d^2$ and $C=a+1$, and induct on $k$.
If $k=1$, an integral curve containing a line is that line.
Thus $L\le\Delta$, which is stronger than the asserted bound.
We can also assume $L>0$ throughout.

Suppose $k\ge2$ and the result holds in dimension $k-1$.
Write $\rho=L/(\Delta q^k)$; by \eqref{eq:point-count},
$0<\rho\le1$. If $\rho<C/q$, then
\[
 L<C\Delta q^{k-1}\le C\Delta q^{\,k-2^{1-k}},
\]
as required. Otherwise put $r=\lceil a/\rho\rceil$.
Then
\[
 r\le\frac a\rho+1\le\frac C\rho\le q.
\]
Proposition~\ref{prop:cut} gives a proper polynomial cut of
degree at most $T=r(q-1)-1<rq$ containing all $L$ lines.

Let $A_1,\ldots,A_m$ be the reduced irreducible components
of this cut. Each has dimension $k-1$, and
\[
 \sum_{i=1}^m\deg A_i\le\Delta T.
\]
Assign each line to one component that contains it. Every
assigned subfamily retains its private points. Put
$\alpha=2^{2-k}$, the saving in dimension $k-1$.
Applying the inductive bound on each component and summing gives
\[
 L\le C\sum_i(\deg A_i)q^{\,k-1-\alpha}
 \le C\Delta Tq^{\,k-1-\alpha}
 \le C\Delta r q^{\,k-\alpha}.
\]
After division by $\Delta q^k$ and use of $r\le C/\rho$,
\[
 \rho\le Crq^{-\alpha}
 \le\frac{C^2}{\rho}q^{-\alpha}.
\]
Hence
\[
 \rho\le Cq^{-\alpha/2}
       =Cq^{-2^{1-k}},
\]
which completes the induction.
\end{proof}

\begin{remark}[Constants and lower-order terms]
The constant was chosen for a short uniform proof and is not optimized.
For the Nikodym problem the conclusion can also be written
\[
 q^d-|N|\le q^{\,d-2^{1-d}+o(1)}.
\]
More explicitly, the $o(1)$ in this display may be taken to be
$\log(8d^2+1)/\log q$. The first exponents are
\[
 \varepsilon_2=\frac12,\quad
 \varepsilon_3=\frac14,\quad
 \varepsilon_4=\frac18,\quad
 \varepsilon_5=\frac1{16}.
\]
No step uses ordinary differentiation, so there is no
large-characteristic hypothesis.
\end{remark}

\end{document}
